\section{Finite limits and colimits}
\label{section-finite-limits}

\noindent
A {\it finite} (co)limit is a (co)limit whose index category is finite,
i.e., the index category has finitely many objects and finitely many
morphisms. A (co)limit is called {\it nonempty} if the index category is
nonempty. A (co)limit is called {\it connected} if the index category is
connected, see
Definition \ref{definition-category-connected}.
It turns out that there are ``enough'' finite index categories.

\begin{lemma}
\label{lemma-finite-diagram-category}
Let $\mathcal{I}$ be a category with
\begin{enumerate}
\item $\Ob(\mathcal{I})$ is finite, and
\item there exist finitely many morphisms
$f_1, \ldots, f_m \in \text{Arrows}(\mathcal{I})$ such
that every morphism of $\mathcal{I}$ is a composition
$f_{j_1} \circ f_{j_2} \circ \ldots \circ f_{j_k}$.
\end{enumerate}
Then there exists a functor $F : \mathcal{J} \to \mathcal{I}$
such that
\begin{enumerate}
\item[(a)] $\mathcal{J}$ is a finite category, and
\item[(b)] for any diagram $M : \mathcal{I} \to \mathcal{C}$ the
(co)limit of $M$ over $\mathcal{I}$ exists if and only if
the (co)limit of $M \circ F$ over $\mathcal{J}$ exists and in this case
the (co)limits are canonically isomorphic.
\end{enumerate}
Moreover, $\mathcal{J}$ is connected (resp.\ nonempty) if and only if
$\mathcal{I}$ is so.
\end{lemma}

\begin{proof}
Say $\Ob(\mathcal{I}) = \{x_1, \ldots, x_n\}$.
Denote $s, t : \{1, \ldots, m\} \to \{1, \ldots, n\}$ the functions
such that $f_j : x_{s(j)} \to x_{t(j)}$.
We set $\Ob(\mathcal{J}) = \{y_1, \ldots, y_n, z_1, \ldots, z_n\}$
Besides the identity morphisms we introduce morphisms
$g_j : y_{s(j)} \to z_{t(j)}$, $j = 1, \ldots, m$ and morphisms
$h_i : y_i \to z_i$, $i = 1, \ldots, n$. Since all of the nonidentity
morphisms in $\mathcal{J}$ go from a $y$ to a $z$ there are no
compositions to define and no associativities to check.
Set $F(y_i) = F(z_i) = x_i$. Set $F(g_j) = f_j$ and $F(h_i) = \text{id}_{x_i}$.
It is clear that $F$ is a functor.
It is clear that $\mathcal{J}$ is finite.
It is clear that $\mathcal{J}$ is connected, resp.\ nonempty
if and only if $\mathcal{I}$ is so.

\medskip\noindent
Let $M : \mathcal{I} \to \mathcal{C}$ be a diagram.
Consider an object $W$ of $\mathcal{C}$ and morphisms
$q_i : W \to M(x_i)$ as in
Definition \ref{definition-limit}.
Then by taking $q_i : W \to M(F(y_i)) = M(F(z_i)) = M(x_i)$ we obtain
a family of maps as in
Definition \ref{definition-limit}
for the diagram $M \circ F$.
Conversely, suppose we are given maps
$qy_i : W \to M(F(y_i))$ and $qz_i : W \to M(F(z_i))$
as in
Definition \ref{definition-limit}
for the diagram $M \circ F$. Since
$$
M(F(h_i)) = \text{id} : M(F(y_i)) = M(x_i) \longrightarrow M(x_i) = M(F(z_i))
$$
we conclude that $qy_i = qz_i$ for all $i$. Set $q_i$ equal to this common
value. The compatibility of
$q_{s(j)} = qy_{s(j)}$ and $q_{t(j)} = qz_{t(j)}$ with the morphism
$M(f_j)$ guarantees that the family $q_i$ is compatible with all morphisms
in $\mathcal{I}$ as by assumption every such morphism is a composition
of the morphisms $f_j$. Thus we have found a canonical bijection
$$
\lim_{B \in \Ob(\mathcal{J})} \Mor_\mathcal{C}(W, M(F(B)))
=
\lim_{A \in \Ob(\mathcal{I})} \Mor_\mathcal{C}(W, M(A))
$$
which implies the statement on limits in the lemma. The statement on colimits
is proved in the same way (proof omitted).
\end{proof}

\begin{lemma}
\label{lemma-fibre-products-equalizers-exist}
Let $\mathcal{C}$ be a category.
The following are equivalent:
\begin{enumerate}
\item Connected finite limits exist in $\mathcal{C}$.
\item Equalizers and fibre products exist in $\mathcal{C}$.
\end{enumerate}
\end{lemma}

\begin{proof}
Since equalizers and fibre products are finite connected
limits we see that (1) implies (2). For the converse, let $\mathcal{I}$
be a finite connected index category. Let
$F : \mathcal{J} \to \mathcal{I}$
be the functor of index categories constructed in the proof of
Lemma \ref{lemma-finite-diagram-category}.
Then we see that we may replace $\mathcal{I}$ by $\mathcal{J}$.
The result is that we may assume that
$\Ob(\mathcal{I}) = \{x_1, \ldots, x_n\} \amalg \{y_1, \ldots, y_m\}$
with $n, m \geq 1$ such that all nonidentity morphisms in $\mathcal{I}$
are morphisms $f : x_i \to y_j$ for some $i$ and $j$.

\medskip\noindent
Suppose that $n > 1$. Since $\mathcal{I}$ is connected there
exist indices $i_1, i_2$ and $j_0$ and morphisms $a : x_{i_1} \to y_{j_0}$
and $b : x_{i_2} \to y_{j_0}$. Consider the category
$$
\mathcal{I}' =
\{x\} \amalg \{x_1, \ldots, \hat x_{i_1}, \ldots, \hat x_{i_2}, \ldots x_n\}
\amalg \{y_1, \ldots, y_m\}
$$
with
$$
\Mor_{\mathcal{I}'}(x, y_j) = \Mor_\mathcal{I}(x_{i_1}, y_j)
\amalg \Mor_\mathcal{I}(x_{i_2}, y_j)
$$
and all other morphism sets the same as in $\mathcal{I}$. For any functor
$M : \mathcal{I} \to \mathcal{C}$ we can construct a functor
$M' : \mathcal{I}' \to \mathcal{C}$ by setting
$$
M'(x) = M(x_{i_1}) \times_{M(a), M(y_{j_0}), M(b)} M(x_{i_2})
$$
and for a morphism $f' : x \to y_j$ corresponding to, say,
$f : x_{i_1} \to y_j$ we set $M'(f) = M(f) \circ \text{pr}_1$.
Then the functor $M$ has a limit if and only if the functor $M'$ has
a limit (proof omitted). Hence by induction we reduce to the case $n = 1$.

\medskip\noindent
If $n = 1$, then the limit of any $M : \mathcal{I} \to \mathcal{C}$ is
the successive equalizer of pairs of maps $x_1 \to y_j$ hence
exists by assumption.
\end{proof}

\begin{lemma}
\label{lemma-almost-finite-limits-exist}
Let $\mathcal{C}$ be a category.
The following are equivalent:
\begin{enumerate}
\item Nonempty finite limits exist in $\mathcal{C}$.
\item Products of pairs and equalizers exist in $\mathcal{C}$.
\item Products of pairs and fibre products exist in $\mathcal{C}$.
\end{enumerate}
\end{lemma}

\begin{proof}
Since products of pairs, fibre products, and equalizers are limits with
nonempty index categories we see that (1) implies both (2) and (3).
Assume (2). Then finite nonempty products and equalizers exist. Hence by
Lemma \ref{lemma-limits-products-equalizers}
we see that finite nonempty limits exist, i.e., (1) holds. Assume (3).
If $a, b : A \to B$ are morphisms of $\mathcal{C}$, then the
equalizer of $a, b$ is
$$
(A \times_{a, B, b} A)\times_{(\text{pr}_1, \text{pr}_2), A \times A, \Delta} A.
$$
Thus (3) implies (2), and the lemma is proved.
\end{proof}

\begin{lemma}
\label{lemma-finite-limits-exist}
Let $\mathcal{C}$ be a category.
The following are equivalent:
\begin{enumerate}
\item Finite limits exist in $\mathcal{C}$.
\item Finite products and equalizers exist.
\item The category has a final object and fibre products exist.
\end{enumerate}
\end{lemma}

\begin{proof}
Since finite products, fibre products, equalizers, and final objects
are limits over finite index categories we see that (1) implies both (2)
and (3). By
Lemma \ref{lemma-limits-products-equalizers}
above we see that (2) implies (1). Assume (3).
Note that the product $A \times B$ is the fibre product over the
final object. If $a, b : A \to B$ are morphisms of $\mathcal{C}$, then the
equalizer of $a, b$ is
$$
(A \times_{a, B, b} A)\times_{(\text{pr}_1, \text{pr}_2), A \times A, \Delta} A.
$$
Thus (3) implies (2) and the lemma is proved.
\end{proof}

\begin{lemma}
\label{lemma-push-outs-coequalizers-exist}
Let $\mathcal{C}$ be a category.
The following are equivalent:
\begin{enumerate}
\item Connected finite colimits exist in $\mathcal{C}$.
\item Coequalizers and pushouts exist in $\mathcal{C}$.
\end{enumerate}
\end{lemma}

\begin{proof}
Omitted. Hint: This is dual to
Lemma \ref{lemma-fibre-products-equalizers-exist}.
\end{proof}

\begin{lemma}
\label{lemma-almost-finite-colimits-exist}
Let $\mathcal{C}$ be a category.
The following are equivalent:
\begin{enumerate}
\item Nonempty finite colimits exist in $\mathcal{C}$.
\item Coproducts of pairs and coequalizers exist in $\mathcal{C}$.
\item Coproducts of pairs and pushouts exist in $\mathcal{C}$.
\end{enumerate}
\end{lemma}

\begin{proof}
Omitted. Hint: This is the dual of
Lemma \ref{lemma-almost-finite-limits-exist}.
\end{proof}

\begin{lemma}
\label{lemma-colimits-exist}
Let $\mathcal{C}$ be a category.
The following are equivalent:
\begin{enumerate}
\item Finite colimits exist in $\mathcal{C}$.
\item Finite coproducts and coequalizers exist in $\mathcal{C}$.
\item The category has an initial object and pushouts exist.
\end{enumerate}
\end{lemma}

\begin{proof}
Omitted. Hint: This is dual to Lemma \ref{lemma-finite-limits-exist}.
\end{proof}





